\chapter{Programação}
\label{cap:programacao-teoria}

\begin{edital}
Desenvolvimento em Java/JavaEE/JakartaEE/JPA/JavaScript; frameworks JUnit,
Hibernate, JSF, Primefaces, Spring, Spring Cloud, Spring Boot; mobile
(Android/iOS); low-code/no-code; análise estática (clean code, SonarQube);
padrões XML, XSLT, UDDI, REST, JSON; DevOps; Git; codificação (transacional,
analítico, mobile, API); reuso; blockchain. (Java em si tem capítulo
próprio: Capítulo \ref{cap:java-teoria}.)
\end{edital}

\section{Ecossistema Spring}

\begin{conceito}[O nó que o Spring desfaz]
Antes do Spring, uma classe de negócio construía suas próprias
dependências com \texttt{new}: o serviço de pedidos decidia, dentro do
próprio código, qual repositório concreto usar. Isso amarrava três coisas
de uma vez --- trocar a implementação exigia editar e recompilar quem a
usava; testar exigia o banco de verdade, porque não havia onde encaixar um
dublê; e a configuração ficava espalhada por todas as classes que
precisavam dela. O ecossistema abaixo existe para desfazer exatamente esse
nó.

\begin{center}
\begin{tabular}{@{}p{2.6cm}p{8.6cm}@{}}
\toprule
\textbf{Framework} & \textbf{Papel} \\
\midrule
Spring & contêiner de inversão de controle (IoC) e injeção de dependência \\
Spring Boot & Spring com autoconfiguração e servidor embutido; elimina
configuração manual e código repetitivo \\
Spring Cloud & ferramentas para sistemas distribuídos/microsserviços
(configuração centralizada, descoberta de serviço, gateway) \\
Hibernate & ORM (mapeamento objeto-relacional), implementação de JPA \\
JUnit & framework de testes de unidade \\
\bottomrule
\end{tabular}
\end{center}
\end{conceito}

\begin{conceito}[IoC e DI: o que o contêiner realmente faz]
\textbf{Inversão de controle (IoC)} é isto: quem decide \textbf{qual}
implementação usar e \textbf{quando} construí-la deixa de ser a própria
classe e passa a ser o \textbf{contêiner}. A classe apenas
\textbf{declara} do que precisa e recebe pronto de fora --- o controle foi
invertido, daí o nome.

\textbf{Injeção de dependência (DI)} é a forma concreta de fazer IoC:
entregar a dependência por \textbf{construtor} (a forma preferida --- a
dependência se torna obrigatória e o objeto já nasce em estado válido),
por \textit{setter}, ou por campo anotado. DI e IoC não são conceitos
rivais: DI é o mecanismo, IoC é o princípio que ele implementa.

O ganho é duplo: \textbf{trocar a implementação} sem tocar em quem a usa,
e \textbf{testar} injetando um dublê no lugar do recurso real (um banco de
dados de teste em vez do de produção, por exemplo). É o
\textit{Dependency Inversion} --- o D do SOLID (Capítulo
\ref{cap:padroes-uml-teoria}) --- executado automaticamente por um
contêiner.

O objeto gerenciado pelo contêiner chama-se \textbf{bean}, e o escopo
padrão do Spring é \textit{singleton} --- \textbf{uma instância por
definição de bean, dentro daquele contêiner específico}. Isso \textbf{não}
é o padrão Singleton do GoF (que faz a própria classe impedir uma segunda
instância via construtor privado) --- é o contêiner quem garante a
instância única, a classe em si nem sabe disso. Mesmo nome, mecanismo e
garantias diferentes.
\end{conceito}

\begin{cuidado}
Spring Boot \textbf{não exige} configuração manual de servidor (ele
embute o Tomcat) nem um servidor de aplicação \textit{standalone}; "Spring
só serve para monólito" é falso. Hibernate é ORM/persistência, não
framework de teste. JUnit testa, não faz mapeamento objeto-relacional.
Spring Cloud endereça sistemas \textbf{distribuídos}; Spring Boot serve a
uma aplicação \textbf{individual}.
\end{cuidado}

\section{Formatos de dados: XML, JSON, XSLT}

\begin{conceito}
\begin{center}
\begin{tabular}{@{}p{2.6cm}p{4.6cm}p{4.6cm}@{}}
\toprule
& \textbf{XML} & \textbf{JSON} \\
\midrule
Estrutura & marcação com tags, verboso & pares chave-valor, leve \\
Uso hoje & legado, arquivo de configuração, SOAP & \textbf{APIs web},
transporte leve de dados \\
\bottomrule
\end{tabular}
\end{center}

\textbf{XSLT}: linguagem que \textbf{transforma XML} em outro formato
(HTML, outro XML, texto puro) --- \textbf{não} transforma JSON, porque
opera sobre a estrutura de árvore específica do XML. \textbf{REST}
tipicamente usa JSON como formato de troca; \textbf{SOAP} usa XML.
\end{conceito}

\begin{cuidado}
XSLT não transforma JSON --- é exclusivo de XML. XML e JSON não são
equivalentes/idênticos: JSON não tem comentários nativos nem um tipo
\texttt{date} nativo (datas em JSON são convenção, geralmente string).
\end{cuidado}

\section{Mobile e low-code}

\begin{conceito}[Nativo, multiplataforma, híbrido]
\begin{center}
\begin{tabular}{@{}p{2.6cm}p{3.4cm}p{4.4cm}@{}}
\toprule
\textbf{Abordagem} & \textbf{O que é} & \textbf{Ferramentas / linguagem} \\
\midrule
Nativo & app específico para cada sistema operacional & Android:
\textbf{Kotlin/Java}; iOS: \textbf{Swift/Objective-C} \\
Multiplataforma & um código-fonte para vários sistemas operacionais &
\textbf{Flutter (Dart)}, \textbf{React Native (JavaScript)}, Xamarin
(.NET), Ionic (web) \\
Híbrido & conteúdo web dentro de um contêiner nativo (\textit{WebView}) &
Ionic, Cordova \\
\bottomrule
\end{tabular}
\end{center}

\textbf{Flutter usa Dart} como linguagem; \textbf{React Native usa
JavaScript}, mas renderiza componentes verdadeiramente nativos (diferente
do \textit{WebView} usado pelo híbrido, que renderiza uma página web
dentro de um contêiner). Nativo entrega melhor desempenho e acesso
completo ao hardware; multiplataforma reduz custo e tempo de
desenvolvimento por manter um único código-fonte.

\textbf{Low-code/no-code}: desenvolvimento visual, com pouca ou nenhuma
escrita manual de código; acelera a entrega e permite que pessoas sem
formação em programação ("\textit{citizen developers}") construam
soluções --- mas limita customização e pode gerar dependência
(\textit{lock-in}) da plataforma escolhida.
\end{conceito}

\begin{cuidado}
Pares que se confundem: Dart pertence ao Flutter, Kotlin ao Android
nativo, Swift ao iOS nativo --- não trocar entre si.
\end{cuidado}

\section{Java EE / web \textit{server-side}: JSF e Primefaces}

\begin{conceito}
\textbf{JSF (\textit{JavaServer Faces})}: framework de UI \textbf{baseado
em componentes}, executado no \textbf{servidor}, orientado a eventos ---
segue o padrão \textbf{MVC}. Tem um ciclo de vida de requisição com fases
bem definidas (restaurar view, aplicar valores, validação, atualizar
modelo, invocar aplicação, renderizar resposta).

\textbf{Primefaces}: biblioteca de \textbf{componentes ricos} construída
\textbf{sobre} JSF (tabelas avançadas, gráficos, Ajax pronto) --- não é um
framework independente, é uma extensão de componentes do JSF.

\textbf{JSP/Servlet}: a camada web mais antiga do Java EE --- o Servlet
processa a requisição, o JSP gera a apresentação (view). JSF é a evolução
baseada em componentes que veio depois.
\end{conceito}

\begin{cuidado}
JSF é \textbf{\textit{server-side} baseado em componentes} --- diferente
de uma SPA \textit{client-side} como React, que roda no navegador.
Primefaces é biblioteca de componentes \textbf{de} JSF, não um framework
à parte.
\end{cuidado}

\section{Versionamento com Git}
\label{sec:git-devops-teoria}

\begin{conceito}[Git é distribuído]
Cada clone do repositório tem o \textbf{histórico completo} --- diferente
do SVN, que é centralizado e lida mal com arquivos binários (cada
alteração exige ida e volta ao servidor central).
\end{conceito}

\begin{conceito}[As três áreas, onde mora quase toda pegadinha de Git]
\begin{center}
diretório de trabalho $\xrightarrow{\texttt{git add}}$ \textbf{stage}
(index) $\xrightarrow{\texttt{git commit}}$ repositório local
$\xrightarrow{\texttt{git push}}$ remoto
\end{center}

O \texttt{commit} grava \textbf{o que está no stage}, não o que está no
diretório de trabalho. Editar um arquivo já versionado e fazer commit
\textbf{sem} \texttt{git add} antes não inclui a alteração no commit --- o
commit sai "vazio" daquela mudança específica, silenciosamente, sem
erro nenhum.
\end{conceito}

\begin{conceito}[Vocabulário de comandos]
\begin{center}
\begin{tabular}{@{}p{2.8cm}p{8.4cm}@{}}
\toprule
\textbf{Comando} & \textbf{O que faz} \\
\midrule
\texttt{git clone} & cria a cópia local inicial de um repositório remoto \\
\texttt{git add} & leva a alteração para o stage \\
\texttt{git commit} & grava no repositório local o que está no stage \\
\texttt{git fetch} & traz referências e objetos do remoto \textbf{sem}
tocar na cópia de trabalho \\
\texttt{git pull} & \texttt{fetch} + integração automática (merge ou
rebase) \\
\texttt{git push} & envia commits locais ao remoto \\
\texttt{git merge} & une duas linhas de desenvolvimento criando um
\textbf{commit de junção} \\
\texttt{git rebase} & \textbf{reaplica} os commits sobre outra base,
produzindo histórico linear \\
\texttt{git revert} & cria um commit \textbf{novo} que anula outro
(seguro em branch já pública/compartilhada) \\
\texttt{git reset --hard} & move a branch e \textbf{descarta} alterações
locais \\
\texttt{git cherry-pick} & copia \textbf{commits escolhidos} para a
branch atual \\
\bottomrule
\end{tabular}
\end{center}

\textbf{Regra de ouro do rebase}: ele reescreve commits (produz novos
identificadores para eles). Não deve ser usado numa branch já
compartilhada com outras pessoas --- é a razão de ser chamado de
"perigoso" nesse contexto: quem já baixou os commits antigos entra em
conflito com os reescritos.

\textbf{Fluxos de trabalho comuns}: branch por funcionalidade + \textit{pull
request}; \textbf{Git Flow} (branches \texttt{main}, \texttt{develop},
\texttt{feature/*}, \texttt{release/*}, \texttt{hotfix/*}, cada uma com
papel definido); \textbf{trunk-based} (integrações curtas e frequentes
diretamente no tronco principal, casa melhor com integração contínua).
\end{conceito}

\begin{cuidado}
\texttt{fetch} \textbf{não} altera a cópia de trabalho (só atualiza
referências locais do remoto); \texttt{pull} altera, porque integra
automaticamente. Se o enunciado disser "\textbf{sem alterar} a cópia de
trabalho", é \texttt{fetch}.
\end{cuidado}

\section{CI, entrega contínua e implantação contínua}

\begin{conceito}[Três práticas, cada uma "terminando" num ponto diferente]
\begin{center}
\begin{tabular}{@{}p{4.4cm}p{6.8cm}@{}}
\toprule
\textbf{Prática} & \textbf{Onde termina} \\
\midrule
\textbf{Integração contínua} (CI) & integra e \textbf{testa} a cada
\textit{push} \\
\textbf{Entrega contínua} (\textit{delivery}) & artefato \textbf{sempre
pronto} para produção; a subida efetiva depende de \textbf{aprovação
manual} \\
\textbf{Implantação contínua} (\textit{deployment}) & o artefato aprovado
vai a produção \textbf{automaticamente}, sem intervenção humana \\
\bottomrule
\end{tabular}
\end{center}

Entrega contínua \textbf{pressupõe} integração contínua --- não dá para
ter um artefato "sempre pronto" sem antes integrar e testar
continuamente. \textbf{DevOps}: cultura de aproximar desenvolvimento e
operação, com automação, medição e responsabilidade compartilhada pelo
que está em produção. \textbf{DevSecOps}: desloca a segurança para a
\textbf{esquerda} (\textit{shift left}) do pipeline --- presente desde o
início do desenvolvimento, não como uma auditoria só no fim.
\end{conceito}

\begin{cuidado}
Se o enunciado disser que a subida a produção \textbf{espera aprovação
humana}, é entrega contínua (\textit{delivery}), não implantação contínua
(\textit{deployment}) --- essa distinção é o par mais trocado do tema.
\end{cuidado}

\section{Contêineres}

\begin{conceito}
\textbf{Docker} constrói (a partir de um \texttt{Dockerfile}) e executa a
\textbf{imagem}, que empacota a aplicação junto com suas dependências.
\textbf{Docker Compose} orquestra vários contêineres em \textbf{um único
host} (uso típico: ambiente de desenvolvimento local). \textbf{Kubernetes}
orquestra contêineres num \textbf{cluster} de várias máquinas: mantém o
número desejado de réplicas, reinicia o que falha, distribui carga entre
elas e realiza atualizações graduais (\textit{rolling update}) --- os
detalhes de arquitetura estão no Capítulo \ref{cap:arquitetura-teoria}.
\end{conceito}

\section{Qualidade de código}

\begin{conceito}
\textbf{Clean Code}: princípios de nomes claros, funções curtas, baixo
acoplamento --- o código deve se explicar por si, sem exigir comentário
extenso. \textbf{SonarQube}: ferramenta de análise \textbf{estática}
(sem executar o código) que detecta \textit{code smells}, bugs em
potencial, vulnerabilidades de segurança, dívida técnica acumulada e mede
cobertura de teste.
\end{conceito}

\begin{conceito}[Onde o resto do edital de Programação está tratado]
\begin{itemize}
\item \textbf{SOLID} e \textbf{orientação a objetos} (herança,
polimorfismo, sobrescrita): Capítulo \ref{cap:java-teoria}.
\item Os \textbf{23 padrões GoF} e o par alta coesão x baixo acoplamento:
Capítulo \ref{cap:padroes-uml-teoria}.
\item \textbf{REST na prática} (métodos HTTP, códigos de resposta,
idempotência): Capítulo \ref{cap:arquitetura-teoria}.
\item \textbf{Reuso, testes e ciclo de vida}: Capítulo
\ref{cap:eng-software-teoria}.
\end{itemize}
\end{conceito}

\section{Blockchain}

\begin{conceito}[O que cada bloco guarda --- e o que ele não guarda]
Uma cadeia de blocos encadeados por \textbf{hash do bloco anterior},
tornando a cadeia \textbf{imutável} na prática (alterar um bloco antigo
mudaria seu hash, quebrando o encadeamento com todos os blocos
seguintes). Cada bloco tipicamente contém: hash do bloco anterior,
\textit{timestamp}, dados/transações, \textit{nonce} (usado no consenso
por prova de trabalho), e a raiz de uma árvore de Merkle (resumo
criptográfico eficiente de todas as transações do bloco).

\textbf{O saldo das carteiras não é armazenado diretamente em nenhum
bloco} --- é \textbf{derivado} do histórico de transações (modelo UTXO
no Bitcoin: soma-se todo o dinheiro recebido por aquele endereço e
subtrai-se o que já foi gasto). É descentralizado, com mecanismo de
consenso (PoW --- \textit{Proof of Work}, ou PoS --- \textit{Proof of
Stake}); Ethereum acrescenta \textbf{contratos inteligentes} (código
executável armazenado e rodado na própria cadeia).
\end{conceito}

\section{Estruturas de dados}
\label{sec:estruturas-dados-teoria}

\begin{conceito}[Cada estrutura é um contrato de custo, não uma definição solta]
A forma mais útil de entender estrutura de dados não é decorar "o que é
uma pilha" --- é reconhecer, diante de um \textbf{requisito}
("inserções e remoções frequentes no meio", "busca por chave em tempo
constante", "listar em ordem e consultar por faixa"), qual estrutura
atende com o menor custo. A escolha da estrutura fixa o custo de cada
operação antes mesmo da primeira linha de código, e é uma escolha cara de
desfazer depois, porque todo o código em volta se apoia nela.

\begin{center}
\begin{tabular}{@{}p{2.9cm}p{4.2cm}p{4.4cm}@{}}
\toprule
\textbf{Estrutura} & \textbf{O que é barato} & \textbf{O que é caro / o
preço} \\
\midrule
Vetor (\textit{array}) & acesso pela posição em $O(1)$ & inserir ou
remover no meio é $O(n)$: os demais elementos \textbf{deslocam} \\
Lista encadeada & inserir/remover em $O(1)$ --- \textbf{desde que você já
esteja no ponto} & chegar ao $i$-ésimo elemento é $O(n)$: precisa
percorrer desde o início \\
Duplamente encadeada & remover um nó já conhecido em $O(1)$, percorrer
nos dois sentidos & dois ponteiros por nó (memória extra) \\
Pilha (LIFO) & empilhar/desempilhar no topo em $O(1)$ & só o topo é
acessível \\
Fila (FIFO) & enfileirar/desenfileirar em $O(1)$ & só as duas pontas são
acessíveis \\
Tabela hash & busca/inserção/remoção em $O(1)$ \textbf{em média} &
\textbf{não guarda ordem}; pior caso $O(n)$ \\
Árvore balanceada & busca em $O(\log n)$ \textbf{mantendo a ordem} & mais
lenta que o hash na busca pontual isolada \\
\bottomrule
\end{tabular}
\end{center}

Duas frases resolvem a maioria das questões do tema: \textbf{vetor troca
inserção por acesso, lista encadeada troca acesso por inserção} --- e o
$O(1)$ da lista pressupõe que você já tem o ponteiro em mãos; se for
preciso \textbf{achar} a posição primeiro, soma-se o $O(n)$ da busca por
cima. E: \textbf{hash e árvore não competem pela mesma vaga} --- se o
enunciado pede \textbf{faixa} de valores ou \textbf{listagem ordenada}, a
resposta é a árvore, mesmo que a tabela hash apareça entre as alternativas
com seu $O(1)$ sedutor (hash não preserva ordem, então não serve a essas
consultas).

\textbf{Busca binária}: $O(\log n)$, mas \textbf{exige a coleção já
ordenada} --- sem essa condição, não se aplica. Em arquivo desordenado e
sem índice, a única saída é a \textbf{busca sequencial}, $O(n)$.
\end{conceito}

\begin{exemplo}[Por que uma tabela hash degrada: fator de carga e colisão]
Uma tabela hash com $m$ posições guardando $n$ chaves tem \textbf{fator de
carga} $\alpha = n/m$. Duas chaves diferentes podem cair na mesma posição
(colisão), tratada de duas formas:

\begin{itemize}
\item \textbf{Encadeamento separado}: cada posição aponta para uma lista.
As chaves moram \textbf{fora} do vetor, então $\alpha$ \textbf{pode passar
de 1} (mais chaves que posições, ficam em listas mais longas); o custo
médio de busca cresce com o comprimento da lista naquela posição.
\item \textbf{Endereçamento aberto}: a chave mora \textbf{dentro} do
próprio vetor, numa posição livre encontrada por \textit{sondagem}. Como
cada posição comporta uma única chave, $\alpha \le 1$ \textbf{sempre} ---
é o que limita o fator de carga a 1 nesse esquema.
\end{itemize}

Caso concreto: endereçamento aberto com \textbf{sondagem linear}
("ocupado? tente a próxima posição") sob ocupação alta, $\alpha \approx
0{,}8$. Só uma posição em cinco está livre, e as ocupadas tendem a se
juntar em \textbf{blocos contíguos}. Uma chave que caia em \textbf{
qualquer} ponto de um bloco precisa caminhar até o fim dele --- e, ao se
acomodar ali, \textbf{aumenta o bloco}, tornando mais provável que a
próxima chave também caia nele. Esse laço de realimentação é o
\textbf{agrupamento primário} (\textit{primary clustering}), e é ele que
empurra a busca do $O(1)$ médio anunciado para perto do $O(n)$ na prática
sob alta ocupação. Remédios: sondagem quadrática ou hash duplo (espalham
a sequência de tentativas de forma menos previsível) e, sobretudo,
\textbf{redimensionar a tabela} quando $\alpha$ ultrapassa um limiar --- é
o $\alpha$, não o número absoluto de chaves, que governa o desempenho.

A mesma lógica reaparece no índice hash de um SGBD (Capítulo
\ref{cap:banco-dados-teoria}): encontra a linha por igualdade em tempo
constante, mas é inútil para \texttt{BETWEEN} ou \texttt{ORDER BY} ---
exatamente porque hash não preserva ordem. Estrutura de dados e índice de
banco são a mesma discussão, vista em dois níveis diferentes.
\end{exemplo}

\begin{cuidado}
Não existe acesso $O(1)$ ao $i$-ésimo elemento de uma \textbf{lista
encadeada} (o $O(1)$ é só da inserção, com o ponteiro já em mãos). Não
existe inserção $O(1)$ no meio de um \textbf{vetor} (há deslocamento de
elementos). Uma tabela hash garante busca em tempo constante \textbf{em
média}, nunca no pior caso absoluto --- desconfie de afirmações que
omitem essa ressalva.
\end{cuidado}

\section{Leitura ativa de código}
\label{sec:leitura-ativa-programacao-teoria}

\begin{conceito}[Protocolo para questões que mostram trecho de código]
Quando a prova mostra um trecho de código pequeno e pergunta a
\textbf{saída} ou o \textbf{efeito}, decorar conceito não resolve --- o
antídoto é \textbf{rodar o código mentalmente}, linha a linha, como um
interpretador, anotando o estado das variáveis num rascunho.

\begin{enumerate}
\item \textbf{Não leia por cima}: releia a linha inteira antes de assumir
o que ela faz --- é comum um enunciado mudar um único parâmetro, operador
ou índice em relação ao que "pareceria óbvio" numa leitura apressada.
\item \textbf{Anote o estado}: para cada variável relevante, escreva o
valor depois de cada linha executada --- não confie na memória depois da
terceira linha de código.
\item \textbf{Desconfie de referência x cópia}: a armadilha mais comum em
qualquer linguagem --- duas variáveis apontando para o \textbf{mesmo}
objeto/array, de modo que mudar uma "muda" a outra também, ou o contrário,
quando a API faz uma cópia de fato.
\item \textbf{Confirme contra a assinatura real} da função/método, não
contra o que pareceria fazer sentido.
\end{enumerate}
\end{conceito}

\begin{conceito}[Quatro fatos de linguagem que decidem a resposta]
\begin{itemize}
\item \textbf{Compilada x interpretada x bytecode}: uma linguagem
\textbf{compilada} traduz para código de máquina \textbf{antes} de rodar
--- um erro de tipo aparece já na compilação. Uma \textbf{interpretada}
executa comando a comando --- o mesmo erro só aparece \textbf{na linha em
que a execução chega}. \textbf{Java} é o meio-termo: compila para
\textbf{bytecode} (\texttt{.class}), e a \textbf{JVM} executa esse
bytecode, compilando os trechos "quentes" para código nativo em tempo de
execução (JIT --- \textit{Just-In-Time}).
\item \textbf{Passagem por valor x por referência}: em Java e em Python,
o que se passa para uma função é sempre o \textbf{valor da referência} ao
objeto. Consequência prática: \textbf{reatribuir} o parâmetro dentro da
função \textbf{não} muda a variável de quem chamou; \textbf{mutar} o
objeto apontado (alterar seu conteúdo interno) muda para os dois lados,
porque ambos apontam para o mesmo objeto. Diante do código, a pergunta
certa não é "é por valor ou por referência?" --- é "esta linha reatribui
ou muda internamente (\textit{muta}) o objeto?".
\item \textbf{Imutável x mutável}: em Python, \texttt{str}, \texttt{int}
e \texttt{tuple} são imutáveis; \texttt{list}, \texttt{dict} e
\texttt{set} são mutáveis. A consequência cobrada: só um objeto imutável
(e por isso \textit{hashável}) serve como \textbf{chave de dicionário} ou
elemento de \texttt{set} --- uma tupla pode, uma lista não. Em Java, a
mesma ideia aparece na \texttt{String} imutável, que é a razão de existir
o \texttt{StringBuilder} (Capítulo \ref{cap:java-teoria}).
\item \textbf{Duck typing}: o que importa é o objeto \textbf{responder}
ao método chamado, não a que classe ele formalmente pertence --- a
compatibilidade é verificada \textbf{no uso}, não na declaração de tipo. É
por isso que Python dispensa a interface explícita que Java exige.
\end{itemize}

\textbf{Aviso de homônimo}: o \textit{decorator} do Python (sintaxe
\texttt{@}) é uma função que recebe outra função e devolve uma versão
embrulhada dela --- é metaprogramação da linguagem. \textbf{Não} é o
padrão Decorator do GoF (Capítulo \ref{cap:padroes-uml-teoria}), que é
composição de objetos em tempo de execução. Mesmo nome, assuntos
completamente diferentes.
\end{conceito}

\begin{exemplo}[NumPy: referência x cópia]
\begin{center}
\begin{tabular}{@{}p{5.4cm}p{5.4cm}@{}}
\toprule
\textbf{Código} & \textbf{Execução mental} \\
\midrule
\texttt{a = np.array([1,2,3])} & \texttt{a} = [1,2,3] (original) \\
\texttt{b = a.view()} & \texttt{b} \textbf{compartilha} o buffer de
\texttt{a}; \texttt{b.base is a} \\
\texttt{a[0] = 99} & muda o buffer compartilhado \\
\texttt{print(b[0])} & imprime \textbf{99} (reflete, pois é \textit{view}) \\
\bottomrule
\end{tabular}
\end{center}

Se a segunda linha fosse \texttt{b = a.copy()}, \texttt{b[0]} continuaria
\textbf{1} mesmo depois de alterar \texttt{a[0]}, porque \texttt{copy()}
rompe o vínculo com o buffer original --- \texttt{b.base} seria
\texttt{None}.
\end{exemplo}

\textbf{Dicionário e a função \texttt{max}}: \texttt{max(dic)} compara as
\textbf{chaves} do dicionário; \texttt{max(dic, key=dic.get)} compara os
\textbf{valores} associados a cada chave, mas devolve a \textbf{chave}
vencedora --- nunca o valor. O parâmetro \texttt{key=} de
\texttt{max}/\texttt{min}/\texttt{sorted} diz \textbf{por onde comparar},
nunca \textbf{o que devolver}.

\textbf{Fatiamento de listas}: \texttt{lista[a:b]} inclui o índice
\texttt{a} e \textbf{exclui} o índice \texttt{b} --- \texttt{nums[1:4]}
devolve \textbf{três} elementos (o tamanho do resultado é \texttt{b - a}).
Índice negativo conta a partir do fim (\texttt{-1} é o último elemento).
